\section{Twisted quotients}\label{sec:descent}

\subsection{Descent of separated checks}\label{sec:descent-lemma}

Keep the setting of \cref{prop:checks}: a family $(h_e)_{e\in G^1}$ in $D(G)$ satisfying (a)--(c) there. Let
$\Gamma'\subseteq\Gamma$ be a subgroup with finite quotient $X'=\Gamma'\backslash G$ and projection $\pi$, and let
$\beta=\sup_e\abs{B(e,R)}$, assumed finite. The \emph{descended relation} on the coordinate set $V'=X'^1$ is
\[
  N'(\pi e)=\pi\bigl(\supp h_e\setminus\{e\}\bigr),
\]
which does not depend on the choice of $e$ in its orbit, by~(a). A word $c\colon V'\to\F$ is \emph{annihilated by the
descended checks} if $\sum_fh_e(f)\,c(\pi f)=0$ for every $e\in G^1$.

\begin{proposition}[Descent]\label{prop:descent}
\formal{AboveHalfRate/Descent.lean}{AboveHalf.CheckSystem.qrel,AboveHalf.CheckSystem.isStorageCode_qrel,AboveHalf.CheckSystem.card_qrel_le,AboveHalf.CheckSystem.card_in_qrel_le,AboveHalf.CheckSystem.le_mul_of_edist_le}
\begin{enumerate}[label=(\alph*)]
\item Every set of words annihilated by the descended checks is a storage code for~$N'$.
\item Both incidence degrees of $N'$ are at most $\beta-1$.
\item If every element of $\Gamma'\setminus\{1\}$ moves every edge by at least $s$, then every coordinate of $N'$ has
paired distance at least $s/R$.
\end{enumerate}
\end{proposition}
\begin{outline}
(a) For $e\in G^1$, split the relation $\sum_fh_e(f)c(\pi f)=0$ into the term $f=e$, which is $c(\pi e)$ because
$h_e(e)=1$, and the other terms, whose coordinates $\pi f$ lie in $N'(\pi e)$. So two annihilated words that agree on
$N'(\pi e)$ agree at $\pi e$.

(b) The recovery set $N'(\pi e)$ is the image of $\supp h_e\setminus\{e\}\subseteq B(e,R)\setminus\{e\}$. For the other
degree fix $p\in G^1$. Every recovery set $N'(\pi e)$ containing $\pi p$ has a representative $e$ with $p\in\supp h_e$
and $e\ne p$: if $\pi y=\pi p$ with $y\in\supp h_{e_0}\setminus\{e_0\}$, translate $e_0$ by the element of $\Gamma'$ that
carries $y$ to $p$. Different recovery sets get different representatives, all in $B(p,R)\setminus\{p\}$.

(c) Give a port $L_f$ the signed class $\operatorname{cls}(f)$ and a port $R_f$ the signed class
$\operatorname{cls}(f)+1$. An incidence edge $L_yR_e$ upstairs has $y\in\supp h_e$ and $y\ne e$, so $y$ and $e$ have
different classes by~(b) of \cref{prop:checks}, and the two ports have the same signed class; they are within distance
$R$ by~(c). Every incidence edge of $N'$ lifts to such an edge at a prescribed lift of either end, as in the proof of~(b).
So a path of length $k$ from $L_{\pi a}$ to $R_{\pi a}$ lifts to a path from $L_a$ to $R_{\gamma a}$ with
$\gamma\in\Gamma'$, along which the signed class is constant and the coordinate moves by at most $kR$. Constancy gives
$\operatorname{cls}(a)=\operatorname{cls}(\gamma a)+1=\operatorname{cls}(a)+\chi(\gamma)+1$, so $\chi(\gamma)=1$ and
$\gamma\ne1$. Hence $s\le\dist(a,\gamma a)\le kR$.
\end{outline}

Part (a) does not use the displacement at all. Neither does (b). Without displacement, however, recovery sets may
contain their own coordinate and paired distances may be short. Part (c) is where the twist pays: a paired path must
travel from a coordinate to a translate of it in the other class, and translates by $\Gamma'$ are far apart.

\begin{lemma}[Validity on quotients]\label{lem:valid-quotient}
\notformalized
Every word of the Tanner code $T(X')$ is annihilated by the descended checks.
\end{lemma}
\begin{informalproof}
For $c\in T(X')$ the function $c\circ\pi$ lies in $T(G)$, and $\sum_fh_e(f)c(\pi f)=\langle h_e,c\circ\pi\rangle=0$ by
\cref{lem:valid}, since $h_e\in D(G)$. The last step is \lean{AboveHalf.pairing_eq_zero_of_mem_iSup}; the Lean library
has no quotient Tanner codes, and its main theorems take the annihilation as a hypothesis.
\end{informalproof}

\subsection{The combinatorial core}\label{sec:descent-core}

\begin{theorem}[Theorem T]\label{thm:T}
\partformal{AboveHalfRate/Main.lean}{AboveHalf.not_Q91At_of_tower,AboveHalf.not_Q91BinaryAt_of_tower}{the passage from a
twisted separated structure to the family of separated checks (\cref{prop:checks}) and the facts about the line-graph
metric used there, and the annihilation of the codes $T(X_j)$ by the descended checks (\cref{lem:valid-quotient}); the
Lean theorems take the family of checks, with its radius and ball sizes, and the annihilation as hypotheses.}
Let $G$ be a locally finite graph and $\Gamma$ a group acting on it freely on vertices and without inversions, with a
$\Gamma$-equivariant local datum over a finite field $\F$ whose local checks interpolate on $t$ coordinates. Suppose $G$
carries a twisted separated structure of degeneracy $t$ and depth $R\ge1$, and let $\beta=\sup_e\abs{B(e,R)}<\infty$.
Let $(\Gamma_j)_{j\ge0}$ be subgroups of $\Gamma$ with finite quotients $X_j=\Gamma_j\backslash G$, such that
$\abs{X_j^1}\to\infty$, the displacement $s(\Gamma_j)\to\infty$, and $\dim_\F T(X_j)\ge\rho\abs{X_j^1}$ for all~$j$.
\begin{enumerate}[label=(\alph*)]
\item The descended relations $N_j$ on $X_j^1$ have incidence degree at most $\beta-1$, every coordinate has paired
distance at least $s(\Gamma_j)/R$, and $T(X_j)$ is an $\F$-linear storage code for $N_j$ of rate at least~$\rho$.
\item If $\rho>1/2$, \cref{prob:q91} has a negative answer at the degree bound $\Delta=\beta-1$, with the fixed alphabet~$\F$.
\item If moreover $\iota\colon\F\to\{0,1\}^b$ is injective and $\rho\log_2\abs\F/b>1/2$, \cref{prob:q91} has a negative answer
for binary codes at the degree bound $b(\beta-1)$.
\end{enumerate}
\end{theorem}
\begin{outline}
\Cref{prop:checks} gives the family of checks. Part (a) is \cref{prop:descent,lem:valid-quotient}; the rate of $T(X_j)$
over $\F$ is $\dim T(X_j)/\abs{X_j^1}$. The paired-distance condition holds in the strongest form: for every radius $K$,
once $s(\Gamma_j)>RK$ no coordinate has paired distance at most~$K$. In (b) the codes $T(X_j)$ violate the conclusion of
\cref{prob:q91}. In (c), the binary blow-up of \cref{sec:binary} multiplies degrees by at most $b$, does not shorten paired
distances, and multiplies rates by $\log_2\abs\F/b$, so its rate is at least $\rho\log_2\abs\F/b$. The Lean statements conclude
$\neg\,$\lean{MergeDepth.Q91At}$(\beta-1)$ and $\neg\,$\lean{MergeDepth.Q91BinaryAt}$(b(\beta-1))$ from a family of
checks satisfying (a)--(c) of \cref{prop:checks}, balls of size at most $\beta$, a tower of subgroups with finite
quotients of size tending to infinity and displacement tending to infinity, and codes on the quotients annihilated by
the descended checks with dimension at least $\rho$ times their length.
\end{outline}

Theorem T separates the two halves of the problem. The combinatorial half asks for an equivariant MDS datum on a tree
with a twisted separated structure. The algebraic half asks for quotients whose Tanner codes are much larger than
counting predicts: dimension counting alone gives $\dim T(X_j)\ge(1-2t/(q+1))\abs{X_j^1}$ for local codes of redundancy
$t$ on a $(q+1)$-regular graph, which is less than $1/2$ whenever $t$ is above a quarter of the degree. The next section
supplies quotients whose Tanner codes have rate at least about $(1-t/(q+1))^2$, the rate of a tensor product of two
local codes.
